\documentclass[10pt,a4paper]{amsart}
\usepackage{amsmath,amssymb,amsthm,graphicx}
\usepackage{stmaryrd}
\usepackage{array}
\title{Disjoint regular collapse-free theories admit modular rule-based generalization algorithms}
\author{Mauricio Ayala-Rincón, David M. Cerna, Temur Kutsia, Christophe Ringeissen}
\date{doi:10.4230/LIPIcs.FSCD.2025.7; CC BY4.0 (official publisher HTML/PDF)}
\begin{document}
\maketitle
\noindent\textit{Editorial note.} Source https://drops.dagstuhl.de/entities/document/10.4230/LIPIcs.FSCD.2025.7. Copyright Mauricio Ayala-Rincón, David M. Cerna, Temur Kutsia and Christophe Ringeissen, CC BY4.0. Theorem19 and complete printed author proof/dependency development are faithfully converted. Supporting Theorem9/Lemma10 and examples are not separately counted.\par\medskip
\noindent\textit{Allowed standard prerequisites (editorial).}\begin{itemize}
\item First-order terms, substitutions, equational generalization and standard equality correspondence for disjoint regular collapse-free theories. Publisher section2 and its cited Lemma1 [46]. Exact equality-correspondence statement and abstraction maps are retained; external standard term/equality prerequisites are not recursively reproved.
\end{itemize}\medskip
% BEGIN S2
 
\par\medskip\noindent\textbf{2 Preliminaries}\par
 
% BEGIN S2.SS1
 
\par\medskip\noindent\textbf{2.1 Terms and Substitutions}\par
 
% BEGIN S2.SS1.p1
 
We consider a first-order alphabet consisting of a signature $\Sigma$ (set of fixed arity function symbols) and a set of variables $\mathcal{V}$. The set of terms $\mathcal{T}(\Sigma,\mathcal{V})$ over $\Sigma$ and $\mathcal{V}$ is defined in the standard way: $t\in\mathcal{T}(\Sigma,\mathcal{V})$ iff $t$ is defined by the grammar $t:=x\mid f(t_{1},\ldots,t_{n})$, where $x\in\mathcal{V}$ and $f\in\Sigma$ is an $n$-ary symbol with $n\geq 0$.\par
 
% END S2.SS1.p1
 
% BEGIN S2.SS1.p2
 
We denote arbitrary function symbols by $f,g,h$, constants by $a,b,c$, variables by $x,y,z,v$, and terms by $s,t,r,u$. The notions of term depth, term size, and a position in a term are defined in the standard way, see, e.g., [6]. By $t|_{p}$, we denote the subterm of $t$ at position $p$, and by $t[s]_{p}$, a term obtained from $t$ by replacing the subterm at position $p$ with the term $s$. For any position $p$ in a term $t$ (including the root position $\epsilon$), $t(p)$ is the symbol at position $p$. The set of all variables in $t$ is denoted by $\mathit{Var}(t)$. A term is called linear if no variable occurs in it more than once, and ground if $\mathit{Var}(t)=\emptyset$.\par
 
% END S2.SS1.p2
 
% BEGIN S2.SS1.p3
 
A sequence of terms $t_{1},\dots,t_{n}$ may be written $\bar{t}$ when $n$ is clear from the context.\par
 
% END S2.SS1.p3
 
% BEGIN S2.SS1.p4
 
Given any signature $\Sigma$ and any finite set of constants $C$ such that $\Sigma\cap C=\emptyset$, the signature $\Sigma\cup C$ is denoted by ${\Sigma}^{C}$.\par
 
% END S2.SS1.p4
 
% BEGIN S2.SS1.p5
 
A substitution is a mapping from $\mathcal{V}$ to $\mathcal{T}(\Sigma,\mathcal{V})$, which is the identity almost everywhere. We use the Greek letters $\sigma,\vartheta,\varphi$ to denote substitutions, except for the identity substitution, which is written as $\mathit{Id}$. We represent substitutions using the usual set notation, e.g., $\sigma=\{x_{1}\mapsto t_{1},\dots,x_{n}\mapsto t_{n}\}$ where $\{x_{1},\dots,x_{n}\}$ is the domain of $\sigma$, denoted by $\mathit{Dom}(\sigma)$. Application of a substitution $\sigma$ to a term $t$, denoted by $t\sigma$, is defined as $x\sigma=\sigma(x)$ and $f(t_{1},\ldots,t_{n})\sigma=f(t_{1}\sigma,\ldots,t_{n}\sigma)$. Substitution composition is defined as a composition of mappings. It is associative but not commutative, with $\mathit{Id}$ playing the role of the unit element. We write $\sigma\vartheta$ for the composition of $\sigma$ with $\vartheta$.\par
 
% END S2.SS1.p5
 
% END S2.SS1
 
% BEGIN S2.SS2
 
\par\medskip\noindent\textbf{2.2 Equational Theories}\par
 
% BEGIN S2.SS2.p1
 
Let $E$ be a set of term pairs, i.e., elements of $\mathcal{T}(\Sigma,\mathcal{V})\times\mathcal{T}(\Sigma,\mathcal{V})$. An equational $\Sigma$-theory generated by $E$, denoted by $=_{E}$, is the least congruence relation on $\mathcal{T}(\Sigma,\mathcal{V})$ that is closed under substitution application and contains $E$. We call $E$ a presentation of $=_{E}$, and the elements of $E$ are called the axioms of $=_{E}$, written as $l=r$. If $t=_{E}u$, we say that $u$ is equal modulo $E$ to $t$. When $E=\emptyset$, every term is equal only to itself, and the theory is called the empty or free theory.\par
 
% END S2.SS2.p1
 
% BEGIN S2.SS2.p2
 
It is common to slightly abuse the terminology in the literature by calling both $E$ and $=_{E}$ an equational theory. For a given $E$, we denote by $\mathsf{sig}(E)$ the set of all function symbols occurring in $E$.\par
 
% END S2.SS2.p2
 
% BEGIN S2.SS2.p3
 
A sequence of $E$-equalities $t_{1}=_{E}u_{1},\dots,t_{n}=_{E}u_{n}$ may be written $\bar{t}=_{E}\bar{u}$ when $n$ is clear from the context.\par
 
% END S2.SS2.p3
 
% BEGIN S2.SS2.p4
 
An axiom $l=r$ is regular if $\mathit{Var}(l)=\mathit{Var}(r)$. An axiom $l=r$ is collapse-free if $l$ and $r$ are non-variable terms. An axiom $l=r$ is shallow if variables can only occur at a position at depth at most $1$ in both $l$ and $r$. An equational theory is regular (resp., collapse-free/shallow) if all its axioms are regular (resp., collapse-free/shallow). An equational theory $E$ is finite if, for each term $t$, there are only finitely many terms $u$ such that $t=_{E}u$. An equational theory $E$ is subterm collapse-free if there are no terms $t,u$ such that $t=_{E}u$ and $u$ is a strict subterm of $t$. A subterm collapse-free theory is necessarily collapse-free, and a finite theory is necessarily regular and subterm collapse-free.\par
 
% END S2.SS2.p4
 
% BEGIN S2.SS2.p5
 
A theory $E$ is syntactic if it has a finite resolvent presentation $RP$, defined as a finite set of axioms $RP$ such that each equality $t=_{E}u$ has an equational proof $t\leftrightarrow_{RP}^{*}u$ with at most one equational step $\leftrightarrow_{RP}$ applied at the root position. Note that any axiom in $RP$ can be applied in both directions: from the left to the right and from the right to the left. Any shallow theory is syntactic [17], and any collapse-free theory with finitary unification is syntactic [29].\par
 
% END S2.SS2.p5
 
% BEGIN S2.SS2.p6
 
One can easily check that $A=\{x*(y*z)=(x*y)*z\}$ (Associativity), $C=\{x*y=y*x\}$ (Commutativity), and $AC=A\cup C$ (Associativity-Commutativity) are regular and collapse-free. Moreover, $A$, $C$ and $AC$ are finite and syntactic [34, 29], but only $C$ is shallow. Since $A$, $C$, and $AC$ are finite theories, they are also subterm collapse-free.\par
 
% END S2.SS2.p6
 
% END S2.SS2
 
% BEGIN S2.SS3
 
\par\medskip\noindent\textbf{2.3 Combination of Equational Theories}\par
 
% BEGIN S2.SS3.p1
 
The rest of the paper assumes that $E_{i}$ is a regular collapse-free $\Sigma_{i}$-theory for $i=1,2$, where $\Sigma_{1}$ and $\Sigma_{2}$ are disjoint signatures. For the sake of brevity, the combined signature $\Sigma_{1}\cup\Sigma_{2}$ is abbreviated to $\Sigma_{1,2}$, and the combined theory $E_{1}\cup E_{2}$ is also denoted by $E_{1,2}$.\par
 
% END S2.SS3.p1
 
% BEGIN S2.SS3.p2
 
A $\Sigma$-rooted term $t$ is a term whose root symbol is in $\Sigma$. Given any ${\Sigma}^{C}_{1,2}$-term $t$ where $t$ is $\Sigma_{i}$-rooted, an alien subterm of $t$ is any subterm $s$ of $t$ which is not $\Sigma_{i}$-rooted and such that any proper superterm of $s$ in $t$ is $\Sigma_{i}$-rooted. The set of alien subterms of $t$ is denoted by $Alien(t)$. Given any term $t$, its height of theory is formally defined as follows: $ht(t)=1+\max\{ht(t^{\prime})\leavevmode\nobreak\ |\leavevmode\nobreak\ t^{\prime}\in Alien(t)\}$. Given a finite set $T$ of ground ${\Sigma}^{C}_{1,2}$-terms, we consider $Aliens=\bigcup_{t\in T}\mathit{Alien}(t)$ and a mapping $\pi:Aliens\rightarrow FC$ such that $FC$ is a set of fresh free constants ($FC\cap C=\emptyset$) and for any $a,a^{\prime}\in Aliens$, $\pi(a)=\pi(a^{\prime})$ iff $a=_{E_{1}\cup E_{2}}a^{\prime}$. Conversely, we define $\pi^{-1}:FC\rightarrow Aliens$ as a mapping such that for any $a\in Aliens$, $\pi^{-1}(\pi(a))=_{E_{1}\cup E_{2}}a$. The $i$-abstraction of any term $t\in T$ is denoted by $t^{\pi_{i}}$ and is inductively defined as follows:\par
  
\par\noindent  $\blacksquare$ 
% BEGIN S2.I1.i1.p1
 
$c^{\pi_{i}}=c$ if $c\in C$,\par
 
% END S2.I1.i1.p1
 
 
\par\noindent  $\blacksquare$ 
% BEGIN S2.I1.i2.p1
 
$(f(t_{1},\dots,t_{n}))^{\pi_{i}}=f(t_{1}^{\pi_{i}},\dots,t_{n}^{\pi_{i}})$ if $f\in\Sigma_{i}$,\par
 
% END S2.I1.i2.p1
 
 
\par\noindent  $\blacksquare$ 
% BEGIN S2.I1.i3.p1
 
$(f(t_{1},\dots,t_{n}))^{\pi_{i}}=\pi(f(t_{1},\dots,t_{n}))$ if $f\in\Sigma_{j}$ for $j\neq i$.\par
 
% END S2.I1.i3.p1
 
  
% END S2.SS3.p2
 
% BEGIN Thmtheorem1
 
\par\medskip\noindent\textbf{Lemma 1 (Correspondence for Equality [46]).}\par
 
% BEGIN Thmtheorem1.p1
 
Let $E_{1}$ and $E_{2}$ be two signature-disjoint regular collapse-free theories. For $i=1,2$ and any ${\Sigma}^{C}_{i}$-rooted terms $t$ and $u$, $t=_{E_{1}\cup E_{2}}u$ iff $t^{\pi_{i}}=_{E_{i}}u^{\pi_{i}}$.\par
 
% END Thmtheorem1.p1
 
% END Thmtheorem1
 
% BEGIN S2.SS3.p3
 
Lemma 1 holds when the component theories $E_{1}$ and $E_{2}$ are regular collapse-free. In the general case where the component theories are arbitrary, we need to assume that $t$ and $u$ are in layer-reduced form [23] to have the same correspondence with their $i$-abstractions.\par
 
% END S2.SS3.p3
 
% END S2.SS3
 
% BEGIN S2.SS4
 
\par\medskip\noindent\textbf{2.4 Generalization Problems}\par
 
% BEGIN S2.SS4.p1
 
Given an equational theory $E$, a term $t$ is more general than $u$ modulo $E$, or is an $E$-generalization of $u$, if there exists a substitution $\sigma$ such that $t\sigma=_{E}u$. In such a case, we write $t\preceq_{E}u$ and also say the term $u$ is less general than $t$ modulo $E$ or that $u$ is an $E$-instance of $t$. The relation $\preceq_{E}$ is a quasi-order and generates the equivalence relation, denoted by $\simeq_{E}$. The strict part of $\preceq_{E}$ is denoted by $\prec_{E}$.\par
 
% END S2.SS4.p1
 
% BEGIN S2.SS4.p2
 
An equational generalization problem is formulated as follows:\par
  Given:  
% BEGIN S2.I2.ix1.p1
 
an equational theory $E$ and two terms $t$ and $u$;\par
 
% END S2.I2.ix1.p1
  Find:  
% BEGIN S2.I2.ix2.p1
 
a term $r$ such that $r\preceq_{E}t$ and $r\preceq_{E}u$ ($r$ is said to be an $E$-generalization of $t$ and $u$).\par
 
% END S2.I2.ix2.p1
   
% END S2.SS4.p2
 
% BEGIN S2.SS4.p3
 
Given an equational theory $E$ and two terms $t$ and $u$, a set of terms $\mathcal{G}$ is called a complete set of $E$-generalizations of $t$ and $u$ if it satisfies the following properties:\par
  
\par\noindent  1. 
% BEGIN S2.I3.i1.p1
 
Soundness: every element of $\mathcal{G}$ is an $E$-generalization of $t$ and $u$,\par
 
% END S2.I3.i1.p1
 
 
\par\noindent  2. 
% BEGIN S2.I3.i2.p1
 
Completeness: for every $E$-generalization $q$ of $t$ and $u$ there exists $r\in\mathcal{G}$ such that $q\preceq_{E}r$.\par
 
% END S2.I3.i2.p1
 
  
The set $\mathcal{G}$ is called a minimal complete set of $E$-generalizations of $t$ and $u$, denoted $\mathit{mcsg}_{E}(t,u)$ if it, in addition, satisfies the following:\par
  
\par\noindent  3. 
% BEGIN S2.I4.i1.p1
 
Minimality: if there exist $r,q\in\mathcal{G}$ such that $r\preceq_{E}q$, then $r=q$.\par
 
% END S2.I4.i1.p1
 
  
% END S2.SS4.p3
 
% BEGIN S2.SS4.p4
 
The existence and cardinality of $\mathit{mcsg}_{E}$ define what is called the generalization type of an equational theory $E$ as follows: $E$ has the generalization type 0 (or is nullary) if there are two terms $t$ and $u$ such that $\mathit{mcsg}_{E}(t,u)$ does not exist. Otherwise, the generalization type of $E$ is\par
  
\par\noindent  $\blacksquare$ 
% BEGIN S2.I5.i1.p1
 
unitary, if $\mathit{mcsg}_{E}(t,u)$ is a singleton for all $t$ and $u$ (in this case, the sole element of $\mathit{mcsg}_{E}(t,u)$ is called a least general $E$-generalization of $t$ and $u$ and is denoted by $\mathit{lgg}_{E}(t,u)$),\par
 
% END S2.I5.i1.p1
 
 
\par\noindent  $\blacksquare$ 
% BEGIN S2.I5.i2.p1
 
finitary, if $\mathit{mcsg}_{E}(t,u)$ is finite for all $t$ and $u$, and there exists at least one pair of terms for which this set is not a singleton,\par
 
% END S2.I5.i2.p1
 
 
\par\noindent  $\blacksquare$ 
% BEGIN S2.I5.i3.p1
 
infinitary, if there exist $t$ and $u$ for which $\mathit{mcsg}_{E}(t,u)$ is infinite.\par
 
% END S2.I5.i3.p1
 
  
% END S2.SS4.p4
 
% BEGIN S2.SS4.p5
 
In the rest of the paper, we consider generalization algorithms defined as terminating inference systems transforming anti-unification triples (AUTs, for short) written $x:t\triangleq u$ where $x$ is a variable, and $t$ and $u$ are the two ground terms to generalize. More precisely, we are focusing on rule-based systems working on configurations $\mathcal{C}$ of the form $A;S;\sigma$, where $A$ and $S$ are sets of AUTs, called the active set and store of $\mathcal{C}$, respectively, and $\sigma$ is a substitution, called the generalization component of $\mathcal{C}$. From now on, $\mathcal{C}$ is said to be a PSS-configuration; it consists of a problem together with a store and a substitution. Given a set of AUTs $P=\bigcup_{i=1}^{n}\{x_{i}:t_{i}\triangleq u_{i}\}$, we associate two substitutions: $\{x_{1}\mapsto t_{1},\allowbreak\dots,x_{n}\mapsto t_{n}\}$ is the left substitution of $P$ and $\{x_{1}\mapsto u_{1},\dots,x_{n}\mapsto u_{n}\}$ is the right substitution of $P$. A set of AUTs may be written using the disjoint union symbol $\uplus$ to isolate a specific AUT from that set.\par
 
% END S2.SS4.p5
 
% END S2.SS4
 
% END S2


% BEGIN S3
 
\par\medskip\noindent\textbf{3 Equational Generalization vs Free Generalization}\par
 
% BEGIN S3.p1
 
The following two theorems characterize equational generalizations with the help of free least general generalizations.\par
 
% END S3.p1
 
% BEGIN Thmtheorem2
 
\par\medskip\noindent\textbf{Theorem 2 (Free $\mathit{lgg}$’s versus $E$-Generalizations).}\par
 
% BEGIN Thmtheorem2.p1
 
Given an equational theory $E$ and two terms $t$ and $u$, if a term $r$ is an $E$-generalization of $t$ and $u$, then there exist terms $t^{\prime}$, $u^{\prime}$, and $r^{\prime}$ such that $t^{\prime}=_{E}t$, $u^{\prime}=_{E}u$, $r^{\prime}$ is a free least general generalization of $t^{\prime}$ and $u^{\prime}$, and $r\preceq_{\emptyset}r^{\prime}$.\par
 
% END Thmtheorem2.p1
 
% END Thmtheorem2
 
% BEGIN S3.SS0.SSS0.P0.SPx1
 
\par\medskip\noindent\textbf{Proof.}\par
 
% BEGIN S3.SS0.SSS0.P0.SPx1.p1
 
From $r$ being an $E$-generalization of $t$ and $u$ we get the existence of two substitutions $\sigma$ and $\vartheta$ such that $r\sigma=_{E}t$ and $r\vartheta=_{E}u$. Take $t^{\prime}=r\sigma$ and $u^{\prime}=r\vartheta$. Then we have $t^{\prime}=_{E}t$ and $u^{\prime}=_{E}u$. Moreover, $r$ is a free generalization of $t^{\prime}$ and $u^{\prime}$. Then there exists a free least general generalization $r^{\prime}$ of $t^{\prime}$ and $u^{\prime}$, and we get $r\preceq_{\emptyset}r^{\prime}$. $\hfill\blacktriangleleft$\par
 
% END S3.SS0.SSS0.P0.SPx1.p1
 
% BEGIN Thmtheorem3
 
\par\medskip\noindent\textbf{Theorem 3 (Completeness and Finiteness through Free $\mathit{lgg}$’s).}\par
 
% BEGIN Thmtheorem3.p1
 
Let $E$ be an equational theory, $t$ and $u$ be two terms, and $\mathcal{G}$ be the set of terms $\mathcal{G}_{E}(t,u)=\{\mathit{lgg}_{\emptyset}(t^{\prime},u^{\prime})\mid t^{\prime}\in[t]_{E},\,u^{\prime}\in[u]_{E}\}$. Then\par
  
\par\noindent  1. 
% BEGIN S3.I1.i1.p1
 
$\mathcal{G}_{E}(t,u)$ is a complete set of $E$-generalizations of $t$ and $u$.\par
 
% END S3.I1.i1.p1
 
 
\par\noindent  2. 
% BEGIN S3.I1.i2.p1
 
$\mathcal{G}_{E}(t,u)$ is not necessarily minimal.\par
 
% END S3.I1.i2.p1
 
 
\par\noindent  3. 
% BEGIN S3.I1.i3.p1
 
If $E$ is a finite theory, then $\mathcal{G}_{E}(t,u)$ is finite modulo variable renaming. Consequently, every finite theory $E$ has the $E$-generalization type at most finitary.\par
 
% END S3.I1.i3.p1
 
 
\par\noindent  4. 
% BEGIN S3.I1.i4.p1
 
For any $E$, if $[t]_{E}$ is finite, then $\mathcal{G}_{E}(t,u)$ (as well as $\mathcal{G}_{E}(u,t)$) is also finite modulo variable renaming.\par
 
% END S3.I1.i4.p1
 
  
% END Thmtheorem3.p1
 
% END Thmtheorem3
 
% END S3.SS0.SSS0.P0.SPx1
 
% BEGIN S3.SS0.SSS0.P0.SPx2
 
\par\medskip\noindent\textbf{Proof.}\par
 
% BEGIN S3.SS0.SSS0.P0.SPx2.p1
  
\par\noindent  1. 
% BEGIN S3.I2.i1.p1
 
Every element $r^{\prime}\in\mathcal{G}$ is an $E$-generalization of $t$ and $u$, because there exist substitutions $\sigma$ and $\vartheta$ such that $r^{\prime}\sigma=t^{\prime}=_{E}t$ and $r^{\prime}\vartheta=u^{\prime}=_{E}u$. Thus, together with Theorem 2, we may conclude that for every $E$-generalization $r$ of $t$ and $u$ there exists an $E$-generalization $r^{\prime}$ of $t$ and $u$ such that $r^{\prime}\in\mathcal{G}$ and $r\preceq_{\emptyset}r^{\prime}$, which implies the completeness of $\mathcal{G}$.\par
 
% END S3.I2.i1.p1
 
 
\par\noindent  2. 
% BEGIN S3.I2.i2.p1
 
Let $E$ be the equational theory that asserts the commutativity of $f$, and take $t=f(a,b)$ and $u=f(a,b)$. Then $\mathcal{G}_{E}(t,u)=\{f(a,b),f(x,y)\}$, because $f(a,b)=\mathit{lgg}_{\emptyset}(f(a,b),f(a,b))$ and $f(x,y)=\mathit{lgg}_{\emptyset}(f(a,b),f(b,a))$. Obviously, $f(x,y)$ is more general than $f(a,b)$ and, hence, $\mathcal{G}_{E}(t,u)$ is not minimal.\par
 
% END S3.I2.i2.p1
 
 
\par\noindent  3. 
% BEGIN S3.I2.i3.p1
 
If $E$ is a finite theory, then $[t]_{E}$ and $[u]_{E}$ are finite sets. Free least general generalizations are unique (modulo variable renaming). Hence, $\mathcal{G}_{E}(t,u)$ is finite modulo renaming. Since $E$-matching is finitary in any finite theory $E$, $\mathcal{G}_{E}(t,u)$ can always be minimized, and we get that $E$-generalization type in finite theories is at most finitary.\par
 
% END S3.I2.i3.p1
 
 
\par\noindent  4. 
% BEGIN S3.I2.i4.p1
 
First, note that the depth of a free generalization of two terms never exceeds the minimum of their depths, and such a generalization does not contain any function symbol that does not appear in both of them. Therefore, for a fixed $t^{\prime}\in[t]_{E}$, the depth of $\mathit{lgg}_{\emptyset}(t^{\prime},u^{\prime})$ is bounded by the depth of $t^{\prime}$, and the set of function symbols that potentially may appear in $\mathit{lgg}_{\emptyset}(t^{\prime},u^{\prime})$ is restricted to the set of function symbols of $t^{\prime}$. Hence, only finitely many terms (over a fixed-arity alphabet like ours) can satisfy these conditions. Hence, whatever $u^{\prime}$ we take, the candidates for $\mathit{lgg}(t^{\prime},u^{\prime})$ may come only from a fixed finite (modulo variable renaming) set that depends on $t^{\prime}$ only: the elements of this set are terms whose free instance is $t^{\prime}$. Therefore, for the given $t^{\prime}$, the set of all terms $\{\mathit{lgg}_{\emptyset}(t^{\prime},u^{\prime})\mid u^{\prime}\in\mathcal{T}(\mathcal{F},\mathcal{V})\}$ is finite. Since by assumption $[t]_{E}$ is finite, we get that $\mathcal{G}_{E}(t,u)$ is a finite union of finite sets, which implies that it is finite.\par
 
% END S3.I2.i4.p1
 
  
$\hfill\blacktriangleleft$\par
 
% END S3.SS0.SSS0.P0.SPx2.p1
 
% BEGIN Thmtheorem4
 
\par\medskip\noindent\textbf{Example 4.}\par
 
% BEGIN Thmtheorem4.p1
 
This example illustrates an infinite $\mathcal{G}_{E}$ for $E$ that is non-finite subterm-collapse free (the simplest class of non-finite theories according to [10]). Let $E=\{f(a,g(x))=f(a,x),f(b,g(x))=f(b,x)\}$, $t=f(a,c)$, and $u=f(b,d)$. Then\par
  
\par  \quad  $\displaystyle\mathcal{G}_{E}(t,u)=\{f(x,y),f(x,g(y)),f(x,g(g(y)),\ldots\}.$\quad  \quad  
  
In fact, this set contains an infinite chain\par
  
\par  \quad  \[f(x,y)\prec_{E}f(x,g(y))\prec_{E}f(x,g(g(y))\prec_{E}\cdots.\]\quad  \quad  
  
% END Thmtheorem4.p1
 
% END Thmtheorem4
 
% END S3.SS0.SSS0.P0.SPx2
 
% END S3


% BEGIN S4
 
\par\medskip\noindent\textbf{4 Combined Generalization in Regular Collapse-Free Theories}\par
 
% BEGIN S4.p1
 
In this section, we consider regular collapse-free theories $E$. When $E$ is regular, the $E$-equivalence class of any ground term only contains ground terms. Thus, any solution of an $E$-matching problem necessarily maps a variable occurring in the problem to a ground term. As a consequence, the following property holds when $E$ is assumed to be regular: given any generalization $r$ of a pair of ground terms $(t,u)$, any substitutions $\phi,\psi$ such that $r\phi=_{E}t$ and $r\psi=_{E}u$ satisfy that $x\phi$ and $x\psi$ are ground for any variable $x\in\mathit{Var}(r)$. This assumption is useful in the following definition, where the considered substitutions are necessarily ground.\par
 
% END S4.p1
 
% BEGIN Thmtheorem5
 
\par\medskip\noindent\textbf{Definition 5 (Generalizations with Solved Instantiations).}\par
 
% BEGIN Thmtheorem5.p1
 
Let $E$ be a regular $\Sigma$-theory. Given any pair of ground ${\Sigma}^{C}$-terms $(t,u)$, an $E$-generalization of $(t,u)$ with respective instantiations $(\phi,\psi)$ is a ${\Sigma}^{C}$-term $r$ together with a pair of substitutions $(\phi,\psi)$ such that $\mathit{Var}(r)=\mathit{Dom}(\phi)=\mathit{Dom}(\psi)$, $r\phi=_{E}t$ and $r\psi=_{E}u$. Given any pair of ground ${\Sigma}^{C}$-terms $(t,u)$, a triple representing an $E$-generalization of $(t,u)$ with solved instantiations (an $E$-GSI triple of $(t,u)$, for short) is any triple $(r,\phi,\psi)$ such that\par
  
\par\noindent  $\blacksquare$ 
% BEGIN S4.I1.i1.p1
 
$r$ is an $E$-generalization of $(t,u)$ with respective instantiations $(\phi,\psi)$,\par
 
% END S4.I1.i1.p1
 
 
\par\noindent  $\blacksquare$ 
% BEGIN S4.I1.i2.p1
 
for any $x\in\mathit{Var}(r)$, there is no non-variable $E$-generalization of $(x\phi,x\psi)$,\par
 
% END S4.I1.i2.p1
 
 
\par\noindent  $\blacksquare$ 
% BEGIN S4.I1.i3.p1
 
for any $x,y\in\mathit{Var}(r)$, $x\phi=_{E}y\phi$ and $x\psi=_{E}y\psi$ implies $x=y$111This third condition may be dropped when considering linear generalizations..\par
 
% END S4.I1.i3.p1
 
  
A set $ST$ of $E$-GSI triples of $(t,u)$ is said to be complete if $\bigcup_{(r,\phi,\psi)\in ST}\{r\}$ is a complete set of $E$-generalizations of $(t,u)$. We say that $E$ admits an algorithm computing a complete set of $E$-generalizations with solved instantiations (GSI algorithm, for short) if there exists a computable function returning, for each input pair of ground ${\Sigma}^{C}$-terms $(t,u)$, a complete set of $E$-GSI triples of $(t,u)$ denoted by $GSI_{E}(t,u)$.\par
 
% END Thmtheorem5.p1
 
% END Thmtheorem5
 
% BEGIN S4.p2
 
A GSI algorithm can be constructed via the repeated application of a generalization algorithm together with a matching algorithm. In the case of a finite theory, the generalization problem is finitary, and the matching problem as well. Consequently, this iteration can be implemented, and it is necessarily terminating in the case of a finite theory since the size of any generalization is necessarily bounded. Thus, we have that:\par
 
% END S4.p2
 
% BEGIN Thmtheorem6
 
\par\medskip\noindent\textbf{Lemma 6.}\par
 
% BEGIN Thmtheorem6.p1
 
Any finite theory admits a GSI algorithm.\par
 
% END Thmtheorem6.p1
 
% END Thmtheorem6
 
% BEGIN Thmtheorem7
 
\par\medskip\noindent\textbf{$\blacktriangleright$ Remark 7.}\par
 
% BEGIN Thmtheorem7.p1
 
In addition to the sketch of the proof given in the paragraph preceding Lemma 6, notice that another proof can be obtained using Theorem 3 and the existence of a GSI algorithm for the empty theory.\par
 
% END Thmtheorem7.p1
 
% END Thmtheorem7
 
% BEGIN Thmtheorem8
 
\par\medskip\noindent\textbf{Example 8.}\par
 
% BEGIN Thmtheorem8.p1
 
In Definition 5, the notion of GSI triple is introduced for any regular theory $E$. In particular, $E$ can be the non-finite regular theory from Example 4. Continuing with that example where $t=f(a,c)$, $u=f(b,d)$, and $\mathcal{G}_{E}(t,u)=\{f(x,y),f(x,g(y)),f(x,g(g(y)),\ldots\}$, one can check that\par
  
\par  \quad  \[GSI_{E}(t,u)=\{(r,\phi,\psi)\mid r\in\mathcal{G}_{E}(t,u),\phi=\{x\mapsto a,y\mapsto c\},\psi=\{x\mapsto b,y\mapsto d\}\}.\]\quad  \quad  
  
% END Thmtheorem8.p1
 
% END Thmtheorem8
 
% BEGIN S4.p3
 
Section 5 introduces another class of theories admitting a GSI algorithm. This class includes non-finite theories.\par
 
% END S4.p3
 
% BEGIN S4.p4
 
In general, for the theories we want to combine, we could imagine deriving GSI algorithms from existing rule-based generalization algorithms transforming PSS-configurations (see Section 6). Then, the resulting GSI algorithms can be combined directly as black-boxes, provided the theories are regular collapse-free.\par
 
% END S4.p4
 
% BEGIN Thmtheorem9
 
\par\medskip\noindent\textbf{Theorem 9 (Modular Property for GSI Algorithm).}\par
 
% BEGIN Thmtheorem9.p1
 
The class of regular collapse-free theories admitting a GSI algorithm and with decidable equality is closed by disjoint union.\par
 
% END Thmtheorem9.p1
 
% END Thmtheorem9
 

% BEGIN S4.SS0.SSS0.P0.SPx1
 
\par\medskip\noindent\textbf{Proof.}\par
 
% BEGIN S4.SS0.SSS0.P0.SPx1.p1
 
The rule-based procedure given in Figure 1 can be shown terminating thanks to a complexity measure defined as follows: to each configuration $(A;S;\sigma)$, we associate a pair $(mht(A),|S|)$ where $mht(A)$ is the multiset of heights of theory $\bigcup_{(x^{\prime}:t^{\prime}\triangleq u^{\prime})\in A}\{ht(t^{\prime}),ht(u^{\prime})\}$, and $|S|$ is the cardinality of $S$. These pairs are ordered lexicographically, the first component being ordered by the multiset extension of the (Noetherian) ordering on positive integers, while the second component is ordered by the (Noetherian) ordering on positive integers. One can easily verify that $E_{i}$-Gen and $E_{1,2}$-Sol strictly decrease the first component of the pair, and $E_{1,2}$-Mer does not increase the first component but strictly decreases the second one.\par
 
% END S4.SS0.SSS0.P0.SPx1.p1
 
% BEGIN S4.SS0.SSS0.P0.SPx1.p2
 
The normal forms with respect to the rule-based procedure are the configurations of the form $(\emptyset;S;\sigma)$ upon which $E_{1,2}$-Mer does not apply. Indeed, any configuration $(A;S;\sigma)$ with $A\neq\emptyset$ is necessarily reducible by either $E_{i}$-Gen or $E_{1,2}$-Sol.\par
 
% END S4.SS0.SSS0.P0.SPx1.p2
 
% BEGIN S4.SS0.SSS0.P0.SPx1.p3
 
For any derivation starting with the input configuration $(\{x:t\triangleq u\},\emptyset,\mathit{Id})$ and leading to a final configuration $(\emptyset;S;\sigma)$, we can associate a tuple $(x\sigma,\phi_{S},\psi_{S})$ where\par
  
\par\noindent  $\blacksquare$ 
% BEGIN S4.I2.i1.p1
 
$\phi_{S}=\{y\mapsto t^{\prime}\leavevmode\nobreak\ |\leavevmode\nobreak\ y:t^{\prime}\triangleq u^{\prime}\in S,y\in\mathit{Var}(x\sigma)\}$,\par
 
% END S4.I2.i1.p1
 
 
\par\noindent  $\blacksquare$ 
% BEGIN S4.I2.i2.p1
 
$\psi_{S}=\{y\mapsto u^{\prime}\leavevmode\nobreak\ |\leavevmode\nobreak\ y:t^{\prime}\triangleq u^{\prime}\in S,y\in\mathit{Var}(x\sigma)\}$.\par
 
% END S4.I2.i2.p1
 
  
Then, the set of all such tuples defines a $GSI_{E_{1}\cup E_{2}}(t,u)$ as stated in Definition 5. This is a consequence of the following facts:\par
  
\par\noindent  $\blacksquare$ 
% BEGIN S4.I3.i1.p1
 
${E_{i}}$-Gen is sound and complete by Lemma 10 and Corollary 11 (see below),\par
 
% END S4.I3.i1.p1
 
 
\par\noindent  $\blacksquare$ 
% BEGIN S4.I3.i2.p1
 
$E_{1,2}$-Sol is sound and complete since $E_{1}$ and $E_{2}$ are collapse-free.\par
 
% END S4.I3.i2.p1
 
  
$\hfill\blacktriangleleft$\par
 
% END S4.SS0.SSS0.P0.SPx1.p3

% BEGIN S4.F1
    
\par
 
$E_{i}$-Gen: Generalization in the $E_{i}$-theory  $\{x:t\triangleq u\}\uplus A;S;\sigma\Longrightarrow(P\setminus S_{i})\cup A;S_{i}\cup S;\sigma\{x\mapsto(r\pi^{-1})\},$\par
     
where\par
     
\par\noindent  $\blacksquare$ 
% BEGIN S4.I4.i1.p1
 
$t(\epsilon)\in{\Sigma}^{C}_{i},u(\epsilon)\in{\Sigma}^{C}_{i}$,\par
 
% END S4.I4.i1.p1
 
 
\par\noindent  $\blacksquare$ 
% BEGIN S4.I4.i2.p1
 
$(r,\phi,\psi)\in GSI_{E_{i}}(t^{\pi_{i}},u^{\pi_{i}})$,\par
 
% END S4.I4.i2.p1
 
 
\par\noindent  $\blacksquare$ 
% BEGIN S4.I4.i3.p1
 
$P=\{y:(y\phi)\pi^{-1}\triangleq(y\psi)\pi^{-1}\leavevmode\nobreak\ |\leavevmode\nobreak\ y\in\mathit{Var}(r)\}$,\par
 
% END S4.I4.i3.p1
 
 
\par\noindent  $\blacksquare$ 
% BEGIN S4.I4.i4.p1
 
$S_{i}=\{y:t^{\prime}\triangleq u^{\prime}\leavevmode\nobreak\ |\leavevmode\nobreak\ (y:t^{\prime}\triangleq u^{\prime})\in P,t^{\prime}(\epsilon)\in{\Sigma}^{C}_{i},u^{\prime}(\epsilon)\in{\Sigma}^{C}_{i}\}$.\par
 
% END S4.I4.i4.p1
 
      
\par
 
$E_{1,2}$-Sol: Solving in the combined theory  $\{x:t\triangleq u\}\uplus A;S;\sigma\Longrightarrow A;S\cup\{x:t\triangleq u\};\sigma,$ if $t(\epsilon)\in\Sigma_{i}$ and $u(\epsilon)\in\Sigma_{3-i}$.\par
      
\par
 
$E_{1,2}$-Mer: Merging in the combined theory  $\emptyset;\{x:t\triangleq u,x^{\prime}:t^{\prime}\triangleq u^{\prime}\}\uplus
S;\sigma\Longrightarrow\emptyset;S\cup\{x:t\triangleq u\};\sigma\{x^{\prime}\mapsto x\},$ if $t=_{E_{1}\cup E_{2}}t^{\prime}$ and $u=_{E_{1}\cup E_{2}}u^{\prime}$.\par
    
Figure 1: Combined Generalization Procedure for Regular Collapse-Free Theories.\par
 
% END S4.F1
 
% BEGIN S4.SS0.SSS0.P0.SPx1.p4
 
We now show the completeness of the algorithm presented in Figure 1 by considering any generalization of the given terms over the combined theory $E_{1}\cup E_{2}$ and showing that we can construct from it an $E_{i}$-generalization, where $i$ depends on the root symbol, that uses the same variable instantiations as the combined generalization. In essence, our combination algorithm builds on the completeness of the algorithms for the individual theories and is thus complete as well.\par
 
% END S4.SS0.SSS0.P0.SPx1.p4
 
% BEGIN S4.SS0.SSS0.P0.SPx1.p5
 
The following technical lemma and corollary are used in the proof of Theorem 9.\par
 
% END S4.SS0.SSS0.P0.SPx1.p5
 
% BEGIN Thmtheorem10
 
\par\medskip\noindent\textbf{Lemma 10 (Correspondence for Generalization).}\par
 
% BEGIN Thmtheorem10.p1
 
Let $E_{1}$ and $E_{2}$ be two signature-disjoint regular collapse-free theories over respectively the signatures $\Sigma_{1}$ and $\Sigma_{2}$, $t$ and $u$ any ${\Sigma}^{C}_{i}$-rooted terms for some $i=1,2$, and $r$ any $E_{1}\cup E_{2}$-generalization $r$ of $(t,u)$ with respective instantiations $(\phi,\psi)$. Consider a bijective mapping $\Pi_{V}:Alien(r)\rightarrow V$ where $V$ is a set of fresh variables and the following terms and substitutions:\par
  
\par\noindent  $\blacksquare$ 
% BEGIN S4.I5.i1.p1
 
$r_{i}$ is the $i$-pure term obtained from $r$ by replacing each $r^{\prime}\in Alien(r)$ by $\Pi_{V}(r^{\prime})$.\par
 
% END S4.I5.i1.p1
 
 
\par\noindent  $\blacksquare$ 
% BEGIN S4.I5.i2.p1
 
$\phi_{i}$ and $\psi_{i}$ are two substitutions such that $\mathit{Dom}(\phi_{i})=\mathit{Dom}(\psi_{i})=\mathit{Var}(r_{i})$ and defined as follows:\par
  
\par\noindent  – 
% BEGIN S4.I5.i2.I1.i1.p1
 
for any $x\in\mathit{Var}(r_{i})\setminus\mathit{Var}(r)$, $x\phi_{i}=((\Pi_{V}^{-1}(x))\phi)^{\pi_{i}}$ and $x\psi_{i}=((\Pi_{V}^{-1}(x))\psi)^{\pi_{i}}$,\par
 
% END S4.I5.i2.I1.i1.p1
 
 
\par\noindent  – 
% BEGIN S4.I5.i2.I1.i2.p1
 
for any $x\in\mathit{Var}(r)$, $x\phi_{i}=(x\phi)^{\pi_{i}}$.\par
 
% END S4.I5.i2.I1.i2.p1
 
  
% END S4.I5.i2.p1
 
  
Then, $r_{i}$ is an $E_{i}$-generalization of $(t^{\pi_{i}},u^{\pi_{i}})$ with respective instantiations $(\phi_{i},\psi_{i})$, and for any $x\in\mathit{Var}(r_{i})\setminus\mathit{Var}(r)$, $\Pi_{V}^{-1}(x)$ is an $E_{1}\cup E_{2}$-generalization of $(x\phi_{i})\pi^{-1},(x\psi_{i})\pi^{-1})$ with respective instantiations $(\phi,\psi)$.\par
 
% END Thmtheorem10.p1
 
% END Thmtheorem10
 
% END S4.SS0.SSS0.P0.SPx1
 

% BEGIN S4.SS0.SSS0.P0.SPx2
 
\par\medskip\noindent\textbf{Proof.}\par
 
% BEGIN S4.SS0.SSS0.P0.SPx2.p1
 
By assumption, $r\phi=_{E_{1}\cup E_{2}}t$ and $r\psi=_{E_{1}\cup E_{2}}u$. By Lemma 1, $(r\phi)^{\pi_{i}}=_{E_{1}\cup E_{2}}t^{\pi_{i}}$ and $(r\psi)^{\pi_{i}}=_{E_{1}\cup E_{2}}u^{\pi_{i}}$. Then, by definition of $r_{i}$, $\phi_{i}$ and $\psi_{i}$, we have $r_{i}\phi_{i}=(r\phi)^{\pi_{i}}$ and $r_{i}\psi_{i}=(r\psi)^{\pi_{i}}$. Consequently, $r_{i}\phi_{i}=_{E_{1}\cup E_{2}}t^{\pi_{i}}$ and $r_{i}\psi_{i}=_{E_{1}\cup E_{2}}u^{\pi_{i}}$.\par
 
% END S4.SS0.SSS0.P0.SPx2.p1
 
% BEGIN S4.SS0.SSS0.P0.SPx2.p2
 
For each $x\in\mathit{Var}(r_{i})\backslash\mathit{Var}(r)$, by definition of $\phi_{i}$ and $\psi_{i}$, we have that $(\Pi_{V}^{-1}(x)\phi)^{\pi_{i}}=x\phi_{i}$ and $(\Pi_{V}^{-1}(x)\psi)^{\pi_{i}}=x\psi_{i}$. Applying $\pi^{-1}$ on these identities, we get $(\Pi_{V}^{-1}(x)\phi)^{\pi_{i}}\pi^{-1}=(x\phi_{i})\pi^{-1}$ and $(\Pi_{V}^{-1}(x)\psi)^{\pi_{i}}\pi^{-1}=(x\psi_{i})\pi^{-1}$. By definition of $\pi^{-1}$, $s=_{E_{1}\cup E_{2}}s^{\pi_{i}}\pi^{-1}$ for any term $s$. Consequently, $(\Pi_{V}^{-1}(x))\phi=_{E_{1}\cup E_{2}}(x\phi_{i})\pi^{-1}$ and $(\Pi_{V}^{-1}(x))\psi=_{E_{1}\cup E_{2}}(x\psi_{i})\pi^{-1}$. $\hfill\blacktriangleleft$\par
 
% END S4.SS0.SSS0.P0.SPx2.p2

% BEGIN Thmtheorem11
 
\par\medskip\noindent\textbf{Corollary 11.}\par
 
% BEGIN Thmtheorem11.p1
 
Let $E_{1}$ and $E_{2}$ be two signature-disjoint regular collapse-free theories over respectively the signatures $\Sigma_{1}$ and $\Sigma_{2}$, and any ${\Sigma}^{C}_{i}$-rooted terms $t,u$ where $i=1,2$. Then, there is no non-variable $E_{1}\cup E_{2}$-generalization of $(t,u)$ iff there is no non-variable $E_{i}$-generalization of $(t^{\pi_{i}},u^{\pi_{i}})$.\par
 
% END Thmtheorem11.p1
 
% END Thmtheorem11
 
% END S4.SS0.SSS0.P0.SPx2
 
% END S4


% BEGIN S5
 
\par\medskip\noindent\textbf{5 Rule-Based Generalization in Finite Syntactic Theories}\par
 
% BEGIN S5.p1
 
In this section, we explore a class of equational theories admitting a rule-based generalization algorithm similar to the one known for the empty theory. For the class of finite syntactic theories, we construct a rule-based algorithm fully parametrized by the axioms of a resolvent presentation of a finite syntactic theory. A similar rule-based algorithm can be found in [27].\par
 
% END S5.p1
 
% BEGIN Thmtheorem12
 
\par\medskip\noindent\textbf{Example 12.}\par
 
% BEGIN Thmtheorem12.p1
 
The class of finite syntactic theories includes $A$, $C$, $AC$, finite shallow theories [17] such as $f(x)=g(x)$, and permutative theories closed by paramodulation [32]. Notice that $\{f(a)=a\}$ and $\{f(f(x))=f(x)\}$ are syntactic but not finite theories.\par
 
% END Thmtheorem12.p1
 
% END Thmtheorem12
 
% BEGIN Thmtheorem13
 
\par\medskip\noindent\textbf{Lemma 13.}\par
 
% BEGIN Thmtheorem13.p1
 
For any finite syntactic theory $E$ with a resolvent presentation $RP$, the inference system given in Figure 4 provides an $E$-generalization algorithm, using the axioms of $RP$ to\par
  
\par\noindent  $\blacksquare$ 
% BEGIN S5.I1.i1.p1
 
check $E$-equality (Figure 2),\par
 
% END S5.I1.i1.p1
 
 
\par\noindent  $\blacksquare$ 
% BEGIN S5.I1.i2.p1
 
solve $E$-matching problems (Figure 3),\par
 
% END S5.I1.i2.p1
 
 
\par\noindent  $\blacksquare$ 
% BEGIN S5.I1.i3.p1
 
control the construction of $E$-generalizations.\par
 
% END S5.I1.i3.p1
 
  
% END Thmtheorem13.p1
 
% END Thmtheorem13
 
% BEGIN S5.SS0.SSS0.P0.SPx1
 
\par\medskip\noindent\textbf{Proof.}\par
 
% BEGIN S5.SS0.SSS0.P0.SPx1.p1
 
The rule-based procedure in Figure 4 terminates since $E$ is assumed to be finite. Indeed, the size of the generalizations associated with configurations is necessarily bounded, and so there are only finitely many applications of Mut. The rule Sol is sound and complete: when an AUT cannot be transformed by Mut or Dec, we cannot have a generalization rooted by a function symbol; a variable is the only possible way to get a generalization for this AUT. The soundness and completeness of Mut and Dec follow because $RP$ is a resolvent presentation of $E$.\par
 
% END S5.SS0.SSS0.P0.SPx1.p1
 
% BEGIN S5.SS0.SSS0.P0.SPx1.p2
 
For any input PSS-configuration $(\{x:t\triangleq u\};\emptyset;\mathit{Id})$ where $t,u$ are ground ${\Sigma}^{C}$-terms, the procedure terminates by computing finitely many PSS-configurations of the form $(\emptyset;S;\sigma)$ and the union of all the terms $x\sigma$ corresponds to a complete set of $E$-generalizations of $t\triangleq u$, since all the rules are sound and complete. $\hfill\blacktriangleleft$\par
 
% END S5.SS0.SSS0.P0.SPx1.p2
 
% BEGIN Thmtheorem14
 
\par\medskip\noindent\textbf{$\blacktriangleright$ Remark 14.}\par
 
% BEGIN Thmtheorem14.p1
 
For the inference system given in Figure 2, an input problem $\{s=t\}$ is such that $s=_{E}t$ iff all the normal forms reached by this inference system are the empty set.\par
 
% END Thmtheorem14.p1
 
% BEGIN Thmtheorem14.p2
 
For the inference system in Figure 3, the set of solutions of an input $E$-matching problem $\{s\leq t\}$ consists of all the normal forms reached by this inference system corresponding to solved forms $\bigcup_{i=1}^{m}\{x_{i}\leq t_{i}\}$ where $x_{1},\dots,x_{m}$ are pairwise distinct variables.\par
 
% END Thmtheorem14.p2
 
% END Thmtheorem14
 
% BEGIN S5.F2
    
\par
 
Mut=:\par
 
\par
 
$\{f(\bar{t})=g(\bar{u})\}\uplus\Gamma\Longrightarrow\{\bar{t}=\bar{l}\theta\}\cup\Gamma$ where $f(\bar{l})=g(\bar{r})$ is a fresh renaming of an axiom in $RP$, and $\theta$ is a substitution such that $\bar{r}\theta=_{RP}\bar{u}$.\par
     
\par
 
Dec=:\par
 
\par
 
$\{f(\bar{t})=f(\bar{u})\}\uplus\Gamma\Longrightarrow\{\bar{t}=\bar{u}\}\cup\Gamma$ where $f\in{\Sigma}^{C}$.\par
    
Figure 2: Procedure for Checking Equality in a finite theory with a resolvent presentation $RP$.\par
 
% END S5.F2
 
% BEGIN S5.F3
    
\par
 
Mut\ensuremath{\leq}: Mutation rule for matching  $\{f(\bar{t})\leq^{?}g(\bar{u})\}\uplus\Gamma\Longrightarrow\{\bar{t}\leq^{?}\bar{l}\theta\}\cup\Gamma$ where $f(\bar{l})=g(\bar{r})$ is a fresh renaming of an axiom in $RP$, and $\theta$ is a substitution such that $\bar{r}\theta=_{RP}\bar{u}$.\par
     
\par
 
Dec\ensuremath{\leq}: Decomposition rule for matching  $\{f(\bar{t})\leq^{?}f(\bar{u})\}\uplus\Gamma\Longrightarrow\{\bar{t}\leq^{?}\bar{u}\}\cup\Gamma$ where $f\in{\Sigma}^{C}$.\par
      
\par
 
Mer\ensuremath{\leq}: Merging rule for matching  $\{x\leq^{?}t,x^{\prime}\leq^{?}t^{\prime}\}\uplus\Gamma\Longrightarrow\{x\leq^{?}t\}\cup\Gamma$ if $t=_{RP}t^{\prime}$.\par
    
Figure 3: Matching Procedure in a finite theory with a resolvent presentation $RP$.\par
 
% END S5.F3
 
% BEGIN S5.F4
    
\par
 
MutLR: Left-Right Mutation rule for generalization  $\{x:f(\bar{t})\triangleq g(\bar{u})\}\uplus A;S;\sigma\Longrightarrow\{y:y\theta\triangleq y\theta^{\prime}\leavevmode\nobreak\ |\leavevmode\nobreak\ y\in\bar{y}\}\cup A;S;\sigma\{x\mapsto h(\bar{y})\}$ where $h(\bar{l})=f(\bar{r})\in RP$, $h(\bar{l}^{\prime})=g(\bar{r}^{\prime})\in RP$, $h(\bar{y})\theta=_{RP}f(\bar{t})$ and $h(\bar{y})\theta^{\prime}=_{RP}g(\bar{u})$.\par
      
\par
 
MutR: Right Mutation rule for generalization  $\{x:f(\bar{t})\triangleq g(\bar{u})\}\uplus A;S;\sigma\Longrightarrow\{y:y\theta\triangleq y\theta^{\prime}\leavevmode\nobreak\ |\leavevmode\nobreak\ y\in\bar{y}\}\cup A;S;\sigma\{x\mapsto f(\bar{y})\}$ where $f(\bar{l}^{\prime})=g(\bar{r}^{\prime})\in RP$, $\theta=\{\bar{y}\mapsto\bar{t}\}$ and $f(\bar{y})\theta^{\prime}=_{RP}g(\bar{u})$.\par
      
\par
 
MutL: Left Mutation rule for generalization  $\{x:f(\bar{t})\triangleq g(\bar{u})\}\uplus A;S;\sigma\Longrightarrow\{y:y\theta\triangleq y\theta^{\prime}\leavevmode\nobreak\ |\leavevmode\nobreak\ y\in\bar{y}\}\cup A;S;\sigma\{x\mapsto g(\bar{y})\}$ where $g(\bar{r}^{\prime})=f(\bar{l}^{\prime})\in RP$, $g(\bar{y})\theta=_{RP}f(\bar{t})$ and $\theta^{\prime}=\{\bar{y}\mapsto\bar{u}\}$.\par
      
\par
 
Dec: Decomposition rule for generalization  $\{x:f(\bar{t})\triangleq f(\bar{u})\}\uplus A;S;\sigma\Longrightarrow\{\bar{y}:\bar{t}=\bar{u}\}\cup A;S;\sigma\{x\mapsto f(\bar{y})\}$ where $f\in{\Sigma}^{C}$ and $\bar{y}$ are fresh variables.\par
     
\par
 
Sol: Solving rule for generalization  $\{x:t\triangleq u\}\uplus A;S;\sigma\Longrightarrow A;S\cup\{x:t\triangleq u\};\sigma$ if no rule in $\{\mathbf{Mut}_{LR},\mathbf{Mut}_{R},\mathbf{Mut}_{L},\mathbf{Dec}\}$ applies.\par
      
\par
 
Mer: Merging rule for generalization  $\emptyset;\{x:t\triangleq u,x^{\prime}:t^{\prime}\triangleq u^{\prime}\}\uplus
S;\sigma\Longrightarrow\emptyset;S\cup\{x:t\triangleq u\};\sigma\{x^{\prime}\mapsto x\}$ if $t=_{RP}t^{\prime}$ and $u=_{RP}u^{\prime}$.\par
    
Figure 4: Generalization Procedure in a finite theory with a resolvent presentation $RP$.\par
 
% END S5.F4
 
% BEGIN S5.SS0.SSS0.P0.SPx1.p3
 
The class of finite syntactic theories is known to be closed by disjoint union [34]. Actually, for a union of two disjoint finite theories with respective resolvent presentations $RP_{1}$ and $RP_{2}$, we have that $RP_{1}\cup RP_{2}$ is a resolvent presentation corresponding to a finite theory. Consequently, the inference system in Figure 4 can be directly applied to any union of disjoint finite syntactic theories for which the resolvent presentations are known. However, finding a resolvent presentation may be a difficult task, especially if $E$-unification is not known to be finitary. In the case we already have a GSI algorithm for a finite theory $E$, we can combine it directly using the approach introduced in Section 4; we don’t have to find a resolvent presentation (if there exists one for $E$).\par
 
% END S5.SS0.SSS0.P0.SPx1.p3
 
% BEGIN Thmtheorem15
 
\par\medskip\noindent\textbf{$\blacktriangleright$ Remark 15.}\par
 
% BEGIN Thmtheorem15.p1
 
The empty theory $\emptyset$ is a finite syntactic theory, and its resolvent presentation is empty. In the case $RP=\emptyset$, one can notice that the inference system from Figure 4 corresponds to the classical rule-based generalization algorithm known for the empty theory.\par
 
% END Thmtheorem15.p1
 
% END Thmtheorem15
 
% END S5.SS0.SSS0.P0.SPx1
 
% END S5


% BEGIN S6
 
\par\medskip\noindent\textbf{6 White-Box Combination of Rule-Based Generalization Algorithms}\par
 
% BEGIN S6.p1
 
This section focuses on the modular construction of generalization algorithms transforming PSS-configurations. Let us now define precisely this kind of rule-based generalization algorithm.\par
 
% END S6.p1
 
% BEGIN Thmtheorem16
 
\par\medskip\noindent\textbf{Definition 16.}\par
 
% BEGIN Thmtheorem16.p1
 
Let $E$ be a regular theory. An $E$-generalization algorithm transforming PSS-configurations is an inference system $\mathfrak{G}$ such that for any initial configuration $(\{x:t\triangleq u\},\emptyset,\mathit{Id})$, it derives a finite set $FC$ of final configurations of the form $(\emptyset;S;\sigma)$ for which $\{x\sigma\leavevmode\nobreak\ |\leavevmode\nobreak\ (\emptyset;S;\sigma)\in FC\}$ is a complete set of $E$-generalizations of $(t,u)$. Furthermore, for any configuration $(A;S;\sigma)$ derived from the initial configuration $(\{x:t\triangleq u\};\emptyset;\mathit{Id})$, the following (invariant) properties must hold:\par
  
\par\noindent  $\blacksquare$ 
% BEGIN S6.I1.i1.p1
 
$x\sigma$ is an $E$-generalization of $(t,u)$ with respective instantiations $(\phi,\psi)$ where $\phi$ (resp., $\psi$) is the left (resp., right) substitution of $A\cup S$,\par
 
% END S6.I1.i1.p1
 
 
\par\noindent  $\blacksquare$ 
% BEGIN S6.I1.i2.p1
 
for any $x^{\prime}:t^{\prime}\triangleq u^{\prime}$ in $S$, there is no non-variable $E$-generalization of $(t^{\prime},u^{\prime})$.\par
 
% END S6.I1.i2.p1
 
  
% END Thmtheorem16.p1
 
% END Thmtheorem16
 
% BEGIN Thmtheorem17
 
\par\medskip\noindent\textbf{$\blacktriangleright$ Remark 17.}\par
 
% BEGIN Thmtheorem17.p1
 
An $E$-generalization algorithm transforming PSS-configurations leads to a GSI-algorithm. Consider an initial configuration $(\{x:t\triangleq u\},\emptyset,\mathit{Id})$ and any triple $(r\sigma,\phi_{S},\psi_{S})$ such that $(\emptyset,S,\sigma)\in FC$ where $FC$ is the set of final configurations given in Definition 16 and $\phi_{S}$ (resp., $\psi_{S}$) denotes the left (resp., right) substitution of $S$. Then, $\{(x\sigma,\phi_{S},\psi_{S})\leavevmode\nobreak\ |\leavevmode\nobreak\ (\emptyset,S,\sigma)\in FC\}$ is actually a $GSI_{E}(t,u)$.\par
 
% END Thmtheorem17.p1
 
% BEGIN Thmtheorem17.p2
 
The inference procedure in Figure 4 corresponds to a generalization algorithm transforming PSS-configurations, thus leading to a GSI algorithm. Conversely, note that a GSI algorithm can be turned into a generalization algorithm transforming PSS-configurations where all derivations are of length $1$.\par
 
% END Thmtheorem17.p2
 
% END Thmtheorem17
 
% BEGIN Thmtheorem18
 
\par\medskip\noindent\textbf{Corollary 18.}\par
 
% BEGIN Thmtheorem18.p1
 
A theory admits a generalization algorithm transforming PSS-configurations iff it admits a GSI algorithm.\par
 
% END Thmtheorem18.p1
 
% END Thmtheorem18
 
% BEGIN S6.p2
 
Generalization algorithms transforming PSS-configurations correspond to GSI algorithms, and so they can be combined using the method given in Figure 1, provided that the respective theories $E_{1}$ and $E_{2}$ are regular collapse-free. Since this method is given as an inference system manipulating PSS-configurations, it can be reworded as an $E_{1}\cup E_{2}$-generalization algorithm transforming PSS-configurations. In this new presentation, we do not directly use the triples from $GSI_{E_{i}}$-sets, but we just apply one step of the inference system corresponding to an $E_{i}$-generalization algorithm transforming PSS-configurations, see Figure 5. This is sufficient to move forward with a solution.\par
 
% END S6.p2
 
% BEGIN S6.p3
 
Given two generalization algorithms $\mathfrak{G}_{1}$ and $\mathfrak{G}_{2}$ transforming PSS-configurations, let $\mathfrak{G}_{1}\oplus\mathfrak{G}_{2}$ be the inference system given by the set of rules $\{\text{{\boldmath$E_{i}$-Step}, {\boldmath$E_{1,2}$-Sol}, {\boldmath$E_{1,2}$-Mer}
}\}$ where $E_{i}$-Step is introduced in Figure 5 while $E_{1,2}$-Sol and $E_{1,2}$-Mer are from Figure 1.\par
 
% END S6.p3
 
% BEGIN S6.F5
  
\par
 
$E_{i}$-Step:\par
 
\par
 
$\{x:t\triangleq u\}\uplus A;S;\sigma\Longrightarrow(A_{i}\pi^{-1})\cup A;(S_{i}\pi^{-1})\cup S;\sigma(\sigma_{i}\pi^{-1})$ where $t(\epsilon)\in{\Sigma}^{C}_{i},u(\epsilon)\in{\Sigma}^{C}_{i}$ and $(\{x:t^{\pi_{i}}\triangleq u^{\pi_{i}}\};\emptyset;\mathit{Id})\longrightarrow_{\mathfrak{G}_{i}}(A_{i};S_{i};\sigma_{i})$.\par
  
Figure 5: Rule calling a generalization algorithm $\mathfrak{G}_{i}$ transforming PSS-configurations.\par
 
% END S6.F5
 

% BEGIN Thmtheorem19
 
\par\medskip\noindent\textbf{Theorem 19 (Modular Construction for Rule-Based Generalization).}\par
 
% BEGIN Thmtheorem19.p1
 
Let $E_{1}$ and $E_{2}$ be two signature-disjoint regular collapse-free theories. If $\mathfrak{G}_{i}$ is an $E_{i}$-generalization algorithm transforming PSS-configurations for $i=1,2$, then $\mathfrak{G}_{1}\oplus\mathfrak{G}_{2}$ is an $E_{1}\cup E_{2}$-generalization algorithm transforming PSS-configurations.\par
 
% END Thmtheorem19.p1
 
% END Thmtheorem19

% BEGIN S6.SS0.SSS0.P0.SPx1
 
\par\medskip\noindent\textbf{Proof.}\par
 
% BEGIN S6.SS0.SSS0.P0.SPx1.p1
 
The inference system $\mathfrak{G}_{1}\oplus\mathfrak{G}_{2}$ is terminating since $\mathfrak{G}_{1}$ and $\mathfrak{G}_{2}$ are terminating, and there are finitely many alternations of calls to $\mathfrak{G}_{1}$ and to $\mathfrak{G}_{2}$ (see Lemma 10). Actually, $\mathfrak{G}_{1}\oplus\mathfrak{G}_{2}$ corresponds to an $E_{1}\cup E_{2}$-generalization algorithm transforming PSS-configurations since the ${E_{i}}$-Step rule from Figure 5 is sound and complete by Lemma 10 and Corollary 11. $\hfill\blacktriangleleft$\par
 
% END S6.SS0.SSS0.P0.SPx1.p1

% BEGIN Thmtheorem20
 
\par\medskip\noindent\textbf{$\blacktriangleright$ Remark 20.}\par
 
% BEGIN Thmtheorem20.p1
 
In both rules $E_{i}$-Step and $E_{i}$-Gen, the constant abstraction $(\cdot)^{\pi_{i}}$ and the term concretization $(\cdot)\pi^{-1}$ should be considered as transparent operations. With $(\cdot)^{\pi_{i}}$, aliens are viewed as free constants, with the restriction that two $E_{1}\cup E_{2}$-equal aliens must be viewed as the same free constant. Conversely, with $(\cdot)\pi^{-1}$, these free constants are viewed back as terms in the combined theory $E_{1}\cup E_{2}$. By slightly abusing notation, if we omit to apply these transformations, then the rules $E_{i}$-Step and $E_{i}$-Gen become easy to express.\par
 
% END Thmtheorem20.p1
 
% END Thmtheorem20
 
% BEGIN S6.SS0.SSS0.P0.SPx1.p2
 
Note that for free, associative, and commutative theories, all four algorithms for the combined theories $(\emptyset)(A)$, $(\emptyset)(C)$, $(A)(C)$, $(\emptyset)(A)(C)$ described in [2] can be obtained as a combination of generalization algorithms transforming PSS-configurations.\par
 
% END S6.SS0.SSS0.P0.SPx1.p2
 
% BEGIN Thmtheorem21
 
\par\medskip\noindent\textbf{Example 21.}\par
 
% BEGIN Thmtheorem21.p1
 
Consider $E_{1}=\{f(x)=g(x)\}$ and $E_{2}=AC(+)$. For $E_{1}$, one can show the existence of an $E_{1}$-generalization algorithm transforming PSS-configurations, say $\mathfrak{G}_{1}$, given by the set of rules $\{\text{$E_{i}$-{Mut}},\ \textbf{Dec},\ \textbf{Sol},\ \textbf{Mer}\}$, where $\textbf{Dec},\textbf{Sol},\textbf{Mer}$ can be found in Figure 4, and the mutation rules $E_{1}$-Mut is defined as follows:\par
 
% END Thmtheorem21.p1
  
% BEGIN Thmtheorem21.p2
 
\par
 
$E_{1}$-Mut:\par
 
\par
 
$\{x:f(t)\triangleq g(u)\}\uplus A;S;\sigma\Longrightarrow\{y:t\triangleq u\}\cup A;S;\sigma\{x\mapsto f(y)\}$\par
 
% END Thmtheorem21.p2
  
% BEGIN Thmtheorem21.p3
 
For $E_{2}$, we call $\mathfrak{G}_{2}$ the $E_{2}$-generalization algorithm transforming PSS-configurations given in Figure 4, with the following resolvent presentation of $AC(+)$ shown in [29]:\par
  
\par  \quad  \[RP=\left\{\begin{array}[]{l}x_{1}+(x_{2}+x_{3})=(x_{1}+x_{2})+x_{3}\\
x_{1}+x_{2}=x_{2}+x_{1}\\
x_{1}+(x_{2}+x_{3})=(x_{1}+x_{3})+x_{2}\\
(x_{1}+x_{2})+x_{3}=(x_{1}+x_{3})+x_{2}\\
x_{1}+(x_{2}+x_{3})=x_{3}+(x_{1}+x_{2})\\
(x_{1}+x_{3})+(x_{2}+x_{4})=(x_{1}+x_{4})+(x_{2}+x_{3})\end{array}\right\}\]\quad  \quad  
  
% END Thmtheorem21.p3
 
% BEGIN Thmtheorem21.p4
 
Let us now detail how the combined generalization algorithm $\mathfrak{G}_{1}\oplus\mathfrak{G}_{2}$ is applied to solve\par
  
\par  \quad  \[x:f(a+(f(a)+b))\triangleq g((a+c)+f(a)).\]\quad  \quad  
  
We focus on two particular derivations leading to final PSS-configurations. The derived PSS-configurations are listed in the table below, where the rightmost column shows the configuration number, while the first column indicates how the configuration in that row is obtained (which generalization algorithm is applied to which configuration):\par
 
% END Thmtheorem21.p4
 
% BEGIN Thmtheorem21.p5
   
\par  \quad  $(\{x:f(a+(f(a)+b))\triangleq g((a+c)+f(a))\};\emptyset;\mathit{Id})$\quad  (0)\quad  
 
\par  $\mathfrak{G}_{1}(0)$\quad  $(\{y_{1}:a+(f(a)+b)\triangleq(a+c)+f(a)\};\emptyset;\sigma_{0}\{x\mapsto f(y_{1})\})$\quad  (1)\quad  
 
\par  $\mathfrak{G}_{2}(1)$\quad  $(\{y_{2}:a\triangleq a+c,y_{3}:f(a)+b\triangleq f(a)\};\emptyset;\sigma_{1}\{y_{1}\mapsto y_{2}+y_{3})\})$\quad  (2)\quad  
 
\par  $\mathfrak{G}_{2}(2)$\quad  $(\{y_{3}:f(a)+b\triangleq f(a)\};\{y_{2}:a\triangleq a+c\};\sigma_{2})$\quad  (3)\quad  
 
\par  $\mathfrak{G}_{2}(3)$\quad  $(\emptyset;\{y_{2}:a\triangleq a+c,y_{3}:f(a)+b\triangleq f(a)\};\sigma_{3})$\quad  (4)\quad  
 
\par  $\mathfrak{G}_{2}(1)$\quad  $(\{y_{4}:f(a)\triangleq f(a),y_{5}:b+a\triangleq a+c\};\emptyset;\sigma_{1}\{y_{1}\mapsto y_{4}+y_{5})\})$\quad  (5)\quad  
 
\par  $\mathfrak{G}_{1}(5)$\quad  $(\{y_{6}:a\triangleq a,y_{5}:b+a\triangleq a+c\};\emptyset;\sigma_{5}\{y_{4}\mapsto f(y_{6})\})$\quad  (6)\quad  
 
\par  $\mathfrak{G}_{1}(6)$\quad  $(\{y_{5}:b+a\triangleq a+c\};\emptyset;\sigma_{6}\{y_{6}\mapsto a\})$\quad  (7)\quad  
 
\par  $\mathfrak{G}_{2}(7)$\quad  $(\{y_{7}:a\triangleq a,y_{8}:b\triangleq c\};\emptyset;\sigma_{9}\{y_{5}\mapsto y_{7}+y_{8}\})$\quad  (8)\quad  
 
\par  $\mathfrak{G}_{1}(8)$\quad  $(\{y_{8}:b\triangleq c\};\emptyset;\sigma_{8}\{y_{7}\mapsto a\})$\quad  (9)\quad  
 
\par  $\mathfrak{G}_{1}(9)$\quad  $(\emptyset;\{y_{8}:b\triangleq c\};\sigma_{9})$\quad  (10)\quad  
   
% END Thmtheorem21.p5
 
% BEGIN Thmtheorem21.p6
 
To get (5), $\textbf{Mut}_{LR}$ is applied with the axioms $x_{1}+(x_{2}+x_{3})=x_{3}+(x_{1}+x_{2})$ and $x_{1}+x_{2}=x_{2}+x_{1}$ in $RP$. To get (8), $\textbf{Mut}_{L}$ is applied with the axiom $x_{1}+x_{2}=x_{2}+x_{1}$ in $RP$. Notice that $\mathfrak{G}_{2}$ could be applied instead of $\mathfrak{G}_{1}$ on (6), (8), (9) since both $\mathfrak{G}_{1}$ and $\mathfrak{G}_{2}$ can handle any AUT between two constants in $C$. Conversely, $\mathfrak{G}_{1}$ could be applied instead of $\mathfrak{G}_{2}$ on (3) since both $\mathfrak{G}_{1}$ and $\mathfrak{G}_{2}$ include a Sol rule. Eventually, (4) and (10) are two final configurations:\par
  
\par\noindent  $\blacksquare$ 
% BEGIN S6.I2.i1.p1
 
$(\emptyset;\{y_{2}:a\triangleq a+c,y_{3}:f(a)+b\triangleq f(a)\};\{x\mapsto f(y_{2}+y_{3}),\dots\}$,\par
 
% END S6.I2.i1.p1
 
 
\par\noindent  $\blacksquare$ 
% BEGIN S6.I2.i2.p1
 
$(\emptyset;\{y_{8}:b\triangleq c\};\{x\mapsto f(f(a)+(a+y_{8})),\dots\})$.\par
 
% END S6.I2.i2.p1
 
  
The resulting generalizations are:\par
  
\par\noindent  $\blacksquare$ 
% BEGIN S6.I3.i1.p1
 
$f(y_{2}+y_{3})$, with respective instantiations $(\{y_{2}\mapsto a,y_{3}\mapsto f(a)+b\},\ \{y_{2}\mapsto a+c,\allowbreak y_{3}\mapsto f(a)\})$,\par
 
% END S6.I3.i1.p1
 
 
\par\noindent  $\blacksquare$ 
% BEGIN S6.I3.i2.p1
 
$f(f(a)+(a+y_{8}))$, with respective instantiations $(\{y_{8}\mapsto b\},\ \{y_{8}\mapsto c\})$.\par
 
% END S6.I3.i2.p1
 
  
% END Thmtheorem21.p6
 
% BEGIN Thmtheorem21.p7
 
The combined generalization algorithm $\mathfrak{G}_{1}\oplus\mathfrak{G}_{2}$ derives additional final configurations. Still, one can check that the corresponding generalizations are always more general than $f(f(a)+(a+y_{8}))$, just like $f(y_{2}+y_{3})$ is more general than $f(f(a)+(a+y_{8}))$.\par
 
% END Thmtheorem21.p7
 
% END Thmtheorem21
 
% END S6.SS0.SSS0.P0.SPx1
 
% END S6

\begin{thebibliography}{99}

% BEGIN bib
 
\par\medskip\noindent\textbf{References}\par
  
\par\noindent  [1]  María Alpuente, Demis Ballis, Angel Cuenca-Ortega, Santiago Escobar, and José Meseguer.  ACUOS${}^{\mbox{2}}$: A High-Performance System for Modular ACU Generalization with Subtyping and Inheritance.  In Francesco Calimeri, Nicola Leone, and Marco Manna, editors, Logics in Artificial Intelligence - 16th European Conference, JELIA 2019, Rende, Italy, May 7-11, 2019, Proceedings, volume 11468 of Lecture Notes in Computer Science, pages 171–181. Springer, 2019.  doi:10.1007/978-3-030-19570-0\_11.  
 
\par\noindent  [2]  María Alpuente, Santiago Escobar, Javier Espert, and José Meseguer.  A modular order-sorted equational generalization algorithm.  Inf. Comput., 235:98–136, 2014.  doi:10.1016/J.IC.2014.01.006.  
 
\par\noindent  [3]  Nino Amiridze and Temur Kutsia.  Anti-unification and natural language processing.  EasyChair Preprint no. 203, EasyChair, 2018.  doi:10.29007/fkrh.  
 
\par\noindent  [4]  Mauricio Ayala-Rincón, David M. Cerna, Andres Felipe Gonzalez Barragan, and Temur Kutsia.  Equational anti-unification over absorption theories.  In Christoph Benzmüller, Marijn J. H. Heule, and Renate A. Schmidt, editors, Automated Reasoning - 12th International Joint Conference, IJCAR 2024, Nancy, France, July 3-6, 2024, Proceedings, Part II, volume 14740 of Lecture Notes in Computer Science, pages 317–337. Springer, 2024.  doi:10.1007/978-3-031-63501-4\_17.  
 
\par\noindent  [5]  Franz Baader.  Unification, weak unification, upper bound, lower bound, and generalization problems.  In Ronald V. Book, editor, Rewriting Techniques and Applications, 4th International Conference, RTA-91, Como, Italy, April 10-12, 1991, Proceedings, volume 488 of Lecture Notes in Computer Science, pages 86–97. Springer, 1991.  doi:10.1007/3-540-53904-2\_88.  
 
\par\noindent  [6]  Franz Baader and Tobias Nipkow.  Term rewriting and all that.  Cambridge University Press, 1998.  doi:10.1017/CBO9781139172752.  
 
\par\noindent  [7]  Franz Baader and Klaus U. Schulz.  Unification in the union of disjoint equational theories: Combining decision procedures.  J. Symb. Comput., 21(2):211–243, 1996.  doi:10.1006/JSCO.1996.0009.  
 
\par\noindent  [8]  Johannes Bader, Andrew Scott, Michael Pradel, and Satish Chandra.  Getafix: learning to fix bugs automatically.  Proc. ACM Program. Lang., 3(OOPSLA):159:1–159:27, 2019.  doi:10.1145/3360585.  
 
\par\noindent  [9]  Adam D. Barwell, Christopher Brown, and Kevin Hammond.  Finding parallel functional pearls: Automatic parallel recursion scheme detection in Haskell functions via anti-unification.  Future Gener. Comput. Syst., 79:669–686, 2018.  doi:10.1016/j.future.2017.07.024.  
 
\par\noindent  [10]  Hans-Jürgen Bürckert, Alexander Herold, and Manfred Schmidt-Schauß.  On equational theories, unification, and (un)decidability.  J. Symb. Comput., 8(1/2):3–49, 1989.  doi:10.1016/S0747-7171(89)80021-5.  
 
\par\noindent  [11]  Jochen Burghardt.  E-generalization using grammars.  Artif. Intell., 165(1):1–35, 2005.  doi:10.1016/J.ARTINT.2005.01.008.  
 
\par\noindent  [12]  David Cao, Rose Kunkel, Chandrakana Nandi, Max Willsey, Zachary Tatlock, and Nadia Polikarpova.  babble: Learning Better Abstractions with E-Graphs and Anti-unification.  Proceedings of the ACM on Programming Languages, 7(POPL):396–424, 2023.  doi:10.1145/3571207.  
 
\par\noindent  [13]  David M. Cerna.  Anti-unification and the theory of semirings.  Theor. Comput. Sci., 848:133–139, 2020.  doi:10.1016/J.TCS.2020.10.020.  
 
\par\noindent  [14]  David M. Cerna and Temur Kutsia.  Idempotent anti-unification.  ACM Trans. Comput. Log., 21(2):10:1–10:32, 2020.  doi:10.1145/3359060.  
 
\par\noindent  [15]  David M. Cerna and Temur Kutsia.  Unital anti-unification: Type and algorithms.  In Zena M. Ariola, editor, 5th International Conference on Formal Structures for Computation and Deduction, FSCD 2020, June 29-July 6, 2020, Paris, France (Virtual Conference), volume 167 of LIPIcs, pages 26:1–26:20. Schloss Dagstuhl – Leibniz-Zentrum für Informatik, 2020.  doi:10.4230/LIPICS.FSCD.2020.26.  
 
\par\noindent  [16]  David M. Cerna and Temur Kutsia.  Anti-unification and generalization: A survey.  In Proceedings of the Thirty-Second International Joint Conference on Artificial Intelligence, IJCAI 2023, 19th-25th August 2023, Macao, SAR, China, pages 6563–6573. ijcai.org, 2023.  doi:10.24963/IJCAI.2023/736.  
 
\par\noindent  [17]  Hubert Comon, Marianne Haberstrau, and Jean-Pierre Jouannaud.  Syntacticness, cycle-syntacticness, and shallow theories.  Inf. Comput., 111(1):154–191, 1994.  doi:10.1006/inco.1994.1043.  
 
\par\noindent  [18]  Andrew Cropper, Sebastijan Dumancic, Richard Evans, and Stephen H. Muggleton.  Inductive logic programming at 30.  Mach. Learn., 111(1):147–172, 2022.  doi:10.1007/S10994-021-06089-1.  
 
\par\noindent  [19]  Reudismam Rolim de Sousa, Gustavo Soares, Rohit Gheyi, Titus Barik, and Loris D’Antoni.  Learning quick fixes from code repositories.  In Cristiano D. Vasconcellos, Karina Girardi Roggia, Vanessa Collere, and Paulo Bousfield, editors, 35th Brazilian Symposium on Software Engineering, SBES 2021, Joinville, Santa Catarina, Brazil, 27 September 2021 - 1 October 2021, pages 74–83. ACM, 2021.  doi:10.1145/3474624.3474650.  
 
\par\noindent  [20]  Eric Domenjoud, Francis Klay, and Christophe Ringeissen.  Combination techniques for non-disjoint equational theories.  In Alan Bundy, editor, Automated Deduction - CADE-12, 12th International Conference on Automated Deduction, Nancy, France, June 26 - July 1, 1994, Proceedings, volume 814 of Lecture Notes in Computer Science, pages 267–281. Springer, 1994.  doi:10.1007/3-540-58156-1\_19.  
 
\par\noindent  [21]  Rui Dong, Zhicheng Huang, Ian Iong Lam, Yan Chen, and Xinyu Wang.  WebRobot: web robotic process automation using interactive programming-by-demonstration.  In Ranjit Jhala and Isil Dillig, editors, PLDI ’22: 43rd ACM SIGPLAN International Conference on Programming Language Design and Implementation, San Diego, CA, USA, June 13 - 17, 2022, pages 152–167. ACM, 2022.  doi:10.1145/3519939.3523711.  
 
\par\noindent  [22]  Serdar Erbatur, Andrew M. Marshall, and Christophe Ringeissen.  Non-disjoint combined unification and closure by equational paramodulation.  In Boris Konev and Giles Reger, editors, Frontiers of Combining Systems - 13th International Symposium, FroCoS 2021, Birmingham, UK, September 8-10, 2021, Proceedings, volume 12941 of Lecture Notes in Computer Science, pages 25–42. Springer, 2021.  doi:10.1007/978-3-030-86205-3\_2.  
 
\par\noindent  [23]  Serdar Erbatur, Andrew M. Marshall, and Christophe Ringeissen.  Combined Hierarchical Matching: the Regular Case.  In Amy P. Felty, editor, 7th International Conference on Formal Structures for Computation and Deduction, FSCD 2022, August 2-5, 2022, Haifa, Israel, volume 228 of LIPIcs, pages 6:1–6:22. Schloss Dagstuhl – Leibniz-Zentrum für Informatik, 2022.  doi:10.4230/LIPICS.FSCD.2022.6.  
 
\par\noindent  [24]  Boris Galitsky.  Developing Enterprise Chatbots - Learning Linguistic Structures.  Springer, 2019.  doi:10.1007/978-3-030-04299-8.  
 
\par\noindent  [25]  Peter Graf.  Substitution tree indexing.  In Jieh Hsiang, editor, Rewriting Techniques and Applications, 6th International Conference, RTA-95, Kaiserslautern, Germany, April 5-7, 1995, Proceedings, volume 914 of Lecture Notes in Computer Science, pages 117–131. Springer, 1995.  doi:10.1007/3-540-59200-8\_52.  
 
\par\noindent  [26]  Alexander Herold.  Combination of unification algorithms in equational theories.  PhD thesis, Kaiserslautern University of Technology, Germany, 1987.  URL: https://d-nb.info/880327979.  
 
\par\noindent  [27]  Jean-Pierre Jouannaud.  Syntactic theories.  In Branislav Rovan, editor, Mathematical Foundations of Computer Science 1990, MFCS’90, Banská Bystrica, Czechoslovakia, August 27-31, 1990, Proceedings, volume 452 of Lecture Notes in Computer Science, pages 15–25. Springer, 1990.  doi:10.1007/BFB0029593.  
 
\par\noindent  [28]  Claude Kirchner.  Méthodes et outils de conception systématique d’algorithmes d’unification dans les théories équationnelles.  Thèses d’état, Université de Nancy I, 1985.  
 
\par\noindent  [29]  Claude Kirchner and Francis Klay.  Syntactic theories and unification.  In Proceedings of the Fifth Annual Symposium on Logic in Computer Science (LICS ’90), Philadelphia, Pennsylvania, USA, June 4-7, 1990, pages 270–277. IEEE Computer Society, 1990.  doi:10.1109/LICS.1990.113753.  
 
\par\noindent  [30]  Hélène Kirchner and Christophe Ringeissen.  Combining symbolic constraint solvers on algebraic domains.  J. Symb. Comput., 18(2):113–155, 1994.  doi:10.1006/JSCO.1994.1040.  
 
\par\noindent  [31]  Boris Konev and Temur Kutsia.  Anti-unification of concepts in description logic EL.  In Chitta Baral, James P. Delgrande, and Frank Wolter, editors, Principles of Knowledge Representation and Reasoning: Proceedings of the Fifteenth International Conference, KR 2016, Cape Town, South Africa, April 25-29, 2016, pages 227–236. AAAI Press, 2016.  URL: http://www.aaai.org/ocs/index.php/KR/KR16/paper/view/12880.  
 
\par\noindent  [32]  Christopher Lynch and Barbara Morawska.  Basic syntactic mutation.  In Andrei Voronkov, editor, Automated Deduction - CADE-18, 18th International Conference on Automated Deduction, Copenhagen, Denmark, July 27-30, 2002, Proceedings, volume 2392 of Lecture Notes in Computer Science, pages 471–485. Springer, 2002.  doi:10.1007/3-540-45620-1\_37.  
 
\par\noindent  [33]  Sonu Mehta, Ranjita Bhagwan, Rahul Kumar, Chetan Bansal, Chandra Shekhar Maddila, B. Ashok, Sumit Asthana, Christian Bird, and Aditya Kumar.  Rex: Preventing bugs and misconfiguration in large services using correlated change analysis.  In Ranjita Bhagwan and George Porter, editors, 17th USENIX Symposium on Networked Systems Design and Implementation, NSDI 2020, Santa Clara, CA, USA, February 25-27, 2020, pages 435–448. USENIX Association, 2020.  URL: https://www.usenix.org/conference/nsdi20/presentation/mehta.  
 
\par\noindent  [34]  Tobias Nipkow.  Proof transformations for equational theories.  In Proceedings of the Fifth Annual Symposium on Logic in Computer Science (LICS ’90), Philadelphia, Pennsylvania, USA, June 4-7, 1990, pages 278–288. IEEE Computer Society, 1990.  doi:10.1109/LICS.1990.113754.  
 
\par\noindent  [35]  Tobias Nipkow.  Combining matching algorithms: The regular case.  J. Symb. Comput., 12(6):633–654, 1991.  doi:10.1016/S0747-7171(08)80145-9.  
 
\par\noindent  [36]  Julian Parsert and Elizabeth Polgreen.  Reinforcement learning and data-generation for syntax-guided synthesis.  In Michael J. Wooldridge, Jennifer G. Dy, and Sriraam Natarajan, editors, Thirty-Eighth AAAI Conference on Artificial Intelligence, AAAI 2024, Thirty-Sixth Conference on Innovative Applications of Artificial Intelligence, IAAI 2024, Fourteenth Symposium on Educational Advances in Artificial Intelligence, EAAI 2014, February 20-27, 2024, Vancouver, Canada, pages 10670–10678. AAAI Press, 2024.  doi:10.1609/AAAI.V38I9.28938.  
 
\par\noindent  [37]  Brigitte Pientka.  Higher-order term indexing using substitution trees.  ACM Trans. Comput. Log., 11(1):6:1–6:40, 2009.  doi:10.1145/1614431.1614437.  
 
\par\noindent  [38]  Gordon D. Plotkin.  A note on inductive generalization.  Machine Intel., 5(1):153–163, 1970.  
 
\par\noindent  [39]  Stanislaw J. Purgal, David M. Cerna, and Cezary Kaliszyk.  Learning higher-order logic programs from failures.  In Luc De Raedt, editor, Proceedings of the Thirty-First International Joint Conference on Artificial Intelligence, IJCAI 2022, Vienna, Austria, 23-29 July 2022, pages 2726–2733. ijcai.org, 2022.  doi:10.24963/IJCAI.2022/378.  
 
\par\noindent  [40]  Mohammad Raza, Sumit Gulwani, and Natasa Milic-Frayling.  Programming by Example Using Least General Generalizations.  In Carla E. Brodley and Peter Stone, editors, Proceedings of the Twenty-Eighth AAAI Conference on Artificial Intelligence, July 27 -31, 2014, Québec City, Québec, Canada, pages 283–290. AAAI Press, 2014.  doi:10.1609/AAAI.V28I1.8744.  
 
\par\noindent  [41]  John C. Reynolds.  Transformational systems and the algebraic structure of atomic formulas.  Machine Intel., 5(1):135–151, 1970.  
 
\par\noindent  [42]  Manfred Schmidt-Schauß.  Unification in a combination of arbitrary disjoint equational theories.  J. Symb. Comput., 8(1/2):51–99, 1989.  doi:10.1016/S0747-7171(89)80022-7.  
 
\par\noindent  [43]  Erik Tidén.  Unification in combinations of collapse-free theories with disjoint sets of function symbols.  In Jörg H. Siekmann, editor, 8th International Conference on Automated Deduction, Oxford, England, July 27 - August 1, 1986, Proceedings, volume 230 of Lecture Notes in Computer Science, pages 431–449. Springer, 1986.  doi:10.1007/3-540-16780-3\_110.  
 
\par\noindent  [44]  Emily Rowan Winter, Vesna Nowack, David Bowes, Steve Counsell, Tracy Hall, Sæmundur Óskar Haraldsson, John R. Woodward, Serkan Kirbas, Etienne Windels, Olayori McBello, Abdurahman Atakishiyev, Kevin Kells, and Matthew W. Pagano.  Towards developer-centered automatic program repair: findings from Bloomberg.  In Abhik Roychoudhury, Cristian Cadar, and Miryung Kim, editors, Proceedings of the 30th ACM Joint European Software Engineering Conference and Symposium on the Foundations of Software Engineering, ESEC/FSE, Singapore, November 14-18, 2022, pages 1578–1588. ACM, 2022.  doi:10.1145/3540250.3558953.  
 
\par\noindent  [45]  Kaiyu Yang and Jia Deng.  Learning symbolic rules for reasoning in quasi-natural language.  Trans. Mach. Learn. Res., 2023, 2023.  URL: https://openreview.net/forum?id=gwRwHUZUgz.  
 
\par\noindent  [46]  Katherine A. Yelick.  Unification in combinations of collapse-free regular theories.  J. Symb. Comput., 3(1/2):153–181, 1987.  doi:10.1016/S0747-7171(87)80025-1.  
 
\par\noindent  [47]  Gonzague Yernaux and Wim Vanhoof.  Anti-unification of unordered goals.  In Florin Manea and Alex Simpson, editors, 30th EACSL Annual Conference on Computer Science Logic, CSL 2022, February 14-19, 2022, Göttingen, Germany (Virtual Conference), volume 216 of LIPIcs, pages 37:1–37:17. Schloss Dagstuhl – Leibniz-Zentrum für Informatik, 2022.  doi:10.4230/LIPICS.CSL.2022.37.  
 
\par\noindent  [48]  Gonzague Yernaux and Wim Vanhoof.  On detecting semantic clones in constraint logic programs.  In 16th IEEE International Workshop on Software Clones, IWSC 2022, Limassol, Cyprus, October 2, 2022, pages 32–38. IEEE, 2022.  doi:10.1109/IWSC55060.2022.00014.  
 
\par\noindent  [49]  Dongjun Youn, Wonho Shin, Jaehyun Lee, Sukyoung Ryu, Joachim Breitner, Philippa Gardner, Sam Lindley, Matija Pretnar, Xiaojia Rao, Conrad Watt, and Andreas Rossberg.  Bringing the WebAssembly Standard up to Speed with SpecTec.  Proc. ACM Program. Lang., 8(PLDI):1559–1584, 2024.  doi:10.1145/3656440.  
  
% END bib
\end{thebibliography}
\end{document}
